\documentclass[10pt,a4paper]{amsart}
\usepackage{amsmath,amssymb,amsthm,graphicx}
% Retained minimal preamble for the Stacks Project chapter more-algebra.tex (More on Algebra).
% The chapter file pinned-original-source.tex carries only \input{preamble} on its first line, so no
% preamble text lives inside the chapter itself.  The Stacks Project uses one single shared preamble for
% every chapter; the snapshot stages it identically as sources/stacks-chow/preamble.tex and
% sources/stacks-etale-cohomology/preamble.tex.  The block below is copied mechanically, line for line,
% from sources/stacks-chow/preamble.tex lines 14-18 (the Xy-pic setup) and 154-217 (the native theorem
% declarations together with the mathematical macros).  Nothing was added, removed or redefined; the
% wrapper class and the amsmath/amssymb/amsthm/graphicx base packages are supplied by the assembler.
% Chapter: more-algebra.tex -- More on Algebra.
% For commutative diagrams we use Xy-pic
\usepackage[all]{xy}

% We use 2cell for 2-commutative diagrams.
\xyoption{2cell}

% Theorem environments.
%
\theoremstyle{plain}
\newtheorem{theorem}[subsection]{Theorem}
\newtheorem{proposition}[subsection]{Proposition}
\newtheorem{lemma}[subsection]{Lemma}

\theoremstyle{definition}
\newtheorem{definition}[subsection]{Definition}
\newtheorem{example}[subsection]{Example}
\newtheorem{exercise}[subsection]{Exercise}
\newtheorem{situation}[subsection]{Situation}

\theoremstyle{remark}
\newtheorem{remark}[subsection]{Remark}
\newtheorem{remarks}[subsection]{Remarks}

\numberwithin{equation}{subsection}

% Macros
%
\def\lim{\mathop{\mathrm{lim}}\nolimits}
\def\colim{\mathop{\mathrm{colim}}\nolimits}
\def\Spec{\mathop{\mathrm{Spec}}}
\def\Hom{\mathop{\mathrm{Hom}}\nolimits}
\def\Ext{\mathop{\mathrm{Ext}}\nolimits}
\def\SheafHom{\mathop{\mathcal{H}\!\mathit{om}}\nolimits}
\def\SheafExt{\mathop{\mathcal{E}\!\mathit{xt}}\nolimits}
\def\Sch{\mathit{Sch}}
\def\Mor{\mathop{\mathrm{Mor}}\nolimits}
\def\Ob{\mathop{\mathrm{Ob}}\nolimits}
\def\Sh{\mathop{\mathit{Sh}}\nolimits}
\def\NL{\mathop{N\!L}\nolimits}
\def\CH{\mathop{\mathrm{CH}}\nolimits}
\def\proetale{{pro\text{-}\acute{e}tale}}
\def\etale{{\acute{e}tale}}
\def\QCoh{\mathit{QCoh}}
\def\Ker{\mathop{\mathrm{Ker}}}
\def\Im{\mathop{\mathrm{Im}}}
\def\Coker{\mathop{\mathrm{Coker}}}
\def\Coim{\mathop{\mathrm{Coim}}}

% Boxtimes
%
\DeclareMathSymbol{\boxtimes}{\mathbin}{AMSa}{"02}

%
% Macros for moduli stacks/spaces
%
\def\QCohstack{\mathcal{QC}\!\mathit{oh}}
\def\Cohstack{\mathcal{C}\!\mathit{oh}}
\def\Spacesstack{\mathcal{S}\!\mathit{paces}}
\def\Quotfunctor{\mathrm{Quot}}
\def\Hilbfunctor{\mathrm{Hilb}}
\def\Curvesstack{\mathcal{C}\!\mathit{urves}}
\def\Polarizedstack{\mathcal{P}\!\mathit{olarized}}
\def\Complexesstack{\mathcal{C}\!\mathit{omplexes}}
% \Pic is the operator that assigns to X its picard group, usage \Pic(X)
% \Picardstack_{X/B} denotes the Picard stack of X over B
% \Picardfunctor_{X/B} denotes the Picard functor of X over B
\def\Pic{\mathop{\mathrm{Pic}}\nolimits}
\def\Picardstack{\mathcal{P}\!\mathit{ic}}
\def\Picardfunctor{\mathrm{Pic}}
\def\Deformationcategory{\mathcal{D}\!\mathit{ef}}

% --- appended 2026-10-12 (Chief): the project own definitions of the reference/slogan/history
% display environments, copied mechanically from the shared Stacks preamble lines 8-12
% (\usepackage{verbatim} plus the three \newenvironment lines, which the chapter preamble keeps
% verbatim from the project preamble). Nothing added, removed or redefined: these environments
% carry no mathematical content and the project itself typesets them as comments. Without these
% three lines a counted statement whose own text uses them fails with Environment slogan undefined.
\usepackage{verbatim}
\newenvironment{reference}{\comment}{\endcomment}
\newenvironment{slogan}{\comment}{\endcomment}
\newenvironment{history}{\comment}{\endcomment}
\usepackage{cleveref}
\crefname{originallemma}{original lemma}{original lemmas}
\Crefname{originallemma}{Original lemma}{Original lemmas}
\makeatletter
\newlabel{categories-definition-dual}{{0FFP}{source}{}{sourcepointer-1}{}}
\newlabel{categories-definition-dual@cref}{{[originallemma][0][]0FFP}{[1][0][]source}{}{}{}}
\makeatother
\crefname{originallemma}{original lemma}{original lemmas}
\Crefname{originallemma}{Original lemma}{Original lemmas}
\makeatletter
\newlabel{lemma-left-dual-module}{{18419}{source}{}{sourcepointer-2}{}}
\newlabel{lemma-left-dual-module@cref}{{[originallemma][0][]18419}{[1][0][]source}{}{}{}}
\makeatother
\title{Left duals in the monoidal category of complexes are degreewise left duals with a sign-twisted differential}
\author{The Stacks Project Authors}
\date{stacks-project:more-algebra.tex:local-snapshot-20261012; GNU FDL 1.2 or later}
\begin{document}
\maketitle
\setcounter{subsection}{1}
\noindent\textit{Editorial note.} All delivered mathematical text is the author's own, copied from the pinned Stacks Project chapter file more-algebra.tex; no statement or proof step was rewritten, reordered or removed.\par\medskip
\noindent\textit{Editorial note.} The chapter file pinned-original-source.tex carries only \textbackslash{}input\{preamble\} on its first line, so no preamble text lives inside the chapter. The retained file sources/stacks-more-algebra/retained-preamble.tex is a mechanical, line-for-line copy of the project's single shared preamble as staged in the snapshot; nothing was added, removed or redefined and no package was added.\par\medskip
\noindent\textit{Editorial note.} The source URL is the source-level Stacks Project URL, as required when more than one card is produced from one source; the per-claim tag is carried by claim\_key.\par\medskip
\noindent\textit{Editorial note.} The counted statement span contains the author's displayed formula for the differential together with its three-part enumerated list and the converse construction, all delivered verbatim. The four references inside the counted proof span are recorded as non-load-bearing navigation pointers.\par\medskip
\noindent\textit{Allowed standard prerequisites (editorial).}\begin{itemize}
\item Dual pairs in a monoidal category and the two triangle identities. invoked at proof lines 18494-18499. Are unwound into the degreewise dual identities of part (3).
\item Left duals of modules and the identification of the dual module with Hom\_R(M, R). invoked at proof lines 18500 and 18509. Gives part (2) and the identification used in the formula for the differential.
\item The sign rule for the differential of a total complex. used at proof lines 18505-18509. Is what produces the factor -(-1)\textasciicircum{}n in the differential of the dual complex.
\end{itemize}\medskip
\ifdefined\hypertarget\hypertarget{sourcepointer-1}{}\fi
\noindent\textit{Source reference.} Original lemma 0FFP: focused-seeds/AGP/sources/stacks-snapshot-20261011/tags.txt, source lines 19294--19294. Cross-chapter navigation pointer into the Stacks Project; the standard fact it supplies is applied inside the counted proof span and the argument resting on it stays there.\par\medskip
\ifdefined\hypertarget\hypertarget{sourcepointer-2}{}\fi
\noindent\textit{Source reference.} Original lemma 18419: sources/stacks-more-algebra/pinned-original-source.tex, source lines 18419--18419. Same-chapter navigation pointer; the statement it names is standard for the section and the argument that consumes it stays inside the counted proof span.\par\medskip
\begin{lemma}
\label{lemma-left-dual-complex}
Let $R$ be a ring. Let $M^\bullet$ be a complex of $R$-modules.
Let $N^\bullet, \eta, \epsilon$ be a left dual of $M^\bullet$
in the monoidal category of complexes of $R$-modules.
Then
\begin{enumerate}
\item $M^\bullet$ and $N^\bullet$ are bounded,
\item $M^n$ and $N^n$ are finite projective $R$-modules,
\item writing $\epsilon = \sum \epsilon_n$
with $\epsilon_n : N^{-n} \otimes_R M^n \to R$ and $\eta = \sum \eta_n$
with $\eta_n : R \to  M^n \otimes_R N^{-n}$ then $(N^{-n}, \eta_n, \epsilon_n)$
is the left dual of $M^n$ as in Lemma \ref{lemma-left-dual-module},
\item the differential $d_N^n : N^n \to N^{n + 1}$ is equal
to $-(-1)^n$ times the map
$$
N^n = \Hom_R(M^{-n}, R)
\xrightarrow{d_M^{-n - 1}}
\Hom_R(M^{-n - 1}, R) = N^{n + 1}
$$
where the equality signs are the identifications from
Lemma \ref{lemma-left-dual-module} part (2).
\end{enumerate}
Conversely, given a bounded complex $M^\bullet$ of finite projective
$R$-modules, setting $N^n = \Hom_R(M^{-n}, R)$ with differentials as above,
setting $\epsilon = \sum \epsilon_n$ with
$\epsilon_n : N^{-n} \otimes_R M^n \to R$ given by evaluation, and
setting $\eta = \sum \eta_n$ with $\eta_n : R \to M^n \otimes_R N^{-n}$
mapping $1$ to $\text{id}_{M_n}$ we obtain a left dual of $M^\bullet$
in the monoidal category of complexes of $R$-modules.
\end{lemma}

\begin{proof}
Since $(1 \otimes \epsilon) \circ (\eta \otimes 1) = \text{id}_{M^\bullet}$
and $(\epsilon \otimes 1) \circ (1 \otimes \eta) = \text{id}_{N^\bullet}$ by
Categories, Definition \ref{categories-definition-dual} we see immediately
that we have
$(1 \otimes \epsilon_n) \circ (\eta_n \otimes 1) = \text{id}_{M^n}$
and $(\epsilon_n \otimes 1) \circ (1 \otimes \eta_n) = \text{id}_{N^{-n}}$
which proves (3). By Lemma \ref{lemma-left-dual-module} we have (2).
Since the sum $\eta = \sum \eta_n$ is finite, we get (1).
Since $\eta = \sum \eta_n$ is a map of complexes
$R \to \text{Tot}(M^\bullet \otimes_R N^\bullet)$ we see that
$$
(d_M^{-n - 1} \otimes 1) \circ \eta_{-n - 1}  +
(-1)^n (1 \otimes d_N^{-n}) \circ \eta_{-n} = 0
$$
by our choice of signs for the differential on
$\text{Tot}(M^\bullet \otimes_R N^\bullet)$.
Unwinding definitions, this proves (4). To see the final statement
of the lemma one reads the above backwards.
\end{proof}

\end{document}
